\section{Problem Set 1: Matrices}

\begin{enumerate}
\item For
\[
A=\begin{pmatrix}
2&-1&4\\
0&3&5
\end{pmatrix},
\]
state the dimensions of \(A\) and identify \(a_{12}\), \(a_{21}\), and \(a_{23}\).

\item Construct the \(3\times 2\) matrix \(B=(b_{ij})\) defined by
\[
b_{ij}=2i-j.
\]

\item Let
\[
A=\begin{pmatrix}1&2\\-3&4\end{pmatrix},
\qquad
B=\begin{pmatrix}5&-1\\2&0\end{pmatrix}.
\]
Compute \(A+B\), \(A-B\), \(3A\), and \(2A-B\).

\item Compute the transpose of
\[
C=\begin{pmatrix}
1&0&-2\\
4&3&5
\end{pmatrix}.
\]
Then state the dimensions of \(C^{\mathsf T}\).

\item Determine which of the following matrices are symmetric:
\[
A=\begin{pmatrix}2&-1\\-1&3\end{pmatrix},
\qquad
B=\begin{pmatrix}0&2\\-2&0\end{pmatrix},
\qquad
C=\begin{pmatrix}1&0\\0&1\end{pmatrix}.
\]

\item Let \(A\) be \(2\times3\), \(B\) be \(3\times4\), and \(C\) be \(4\times2\). Determine which of the products \(AB\), \(BA\), \(BC\), \(CB\), \(AC\), and \(CA\) are defined. State the dimensions of every defined product.

\item Compute
\[
AB
\quad\text{for}\quad
A=\begin{pmatrix}1&2&-1\\0&3&4\end{pmatrix},
\qquad
B=\begin{pmatrix}2&1\\-1&0\\3&-2\end{pmatrix}.
\]

\item For
\[
A=\begin{pmatrix}1&2\\0&1\end{pmatrix},
\qquad
B=\begin{pmatrix}2&0\\1&3\end{pmatrix},
\]
compute both \(AB\) and \(BA\). Explain what the result shows about matrix multiplication.

\item Let
\[
A=\begin{pmatrix}2&-1\\3&4\end{pmatrix},
\qquad
x=\begin{pmatrix}5\\-2\end{pmatrix}.
\]
Compute \(Ax\) and interpret the result as a new column vector.

\item Verify directly that
\[
AI_2=I_2A=A
\]
for
\[
A=\begin{pmatrix}3&-2\\1&5\end{pmatrix}.
\]

\item Let
\[
D=\operatorname{diag}(2,-1,3).
\]
Compute \(D^2\) and \(D^3\). Describe the pattern.

\item Starting with
\[
A=\begin{pmatrix}
1&2&-1\\
3&0&4\\
-2&1&5
\end{pmatrix},
\]
perform the row operation \(R_2\leftarrow R_2-3R_1\).

\item Starting with
\[
A=\begin{pmatrix}
1&2\\
3&4
\end{pmatrix},
\]
first interchange \(R_1\) and \(R_2\), and then perform \(R_2\leftarrow R_2-\frac13R_1\). Write the matrix after each step.

\item Find \(x,y,z,w\) if
\[
\begin{pmatrix}
x&2y\\
z-1&w+3
\end{pmatrix}
=
\begin{pmatrix}
4&6\\
2&1
\end{pmatrix}.
\]

\item Find the matrix \(X\) satisfying
\[
2X+\begin{pmatrix}1&-1\\0&3\end{pmatrix}
=
\begin{pmatrix}5&3\\4&-1\end{pmatrix}.
\]

\item Verify the distributive law \(A(B+C)=AB+AC\) for
\[
A=\begin{pmatrix}1&2\\-1&0\end{pmatrix},
\quad
B=\begin{pmatrix}2&1\\0&3\end{pmatrix},
\quad
C=\begin{pmatrix}-1&4\\2&0\end{pmatrix}.
\]

\item The columns of
\[
A=\begin{pmatrix}
1&0\\
2&1\\
-1&3
\end{pmatrix}
\]
are vectors \(u\) and \(v\). Write \(u\) and \(v\), and express
\[
A\begin{pmatrix}2\\-1\end{pmatrix}
\]
as a linear combination of these columns.

\item The matrix
\[
T=\begin{pmatrix}2&0\\0&-1\end{pmatrix}
\]
acts on points in the plane by matrix multiplication. Find the images of
\[
(1,2)^{\mathsf T},\quad (-3,1)^{\mathsf T},\quad (0,4)^{\mathsf T}.
\]
Describe what \(T\) does geometrically.

\item A shop records sales of two products on three days in the matrix
\[
S=\begin{pmatrix}
12&8\\
15&11\\
9&14
\end{pmatrix}.
\]
Explain what the rows and columns could represent. Compute the total sales of each product using matrix or vector notation.

\item For each expression below, decide whether it is defined. If it is defined, compute it; if it is not, explain the dimension mismatch:
\[
A=\begin{pmatrix}1&2&3\\0&1&-1\end{pmatrix},
\quad
B=\begin{pmatrix}2&0\\1&4\\-2&3\end{pmatrix},
\quad
u=\begin{pmatrix}1\\-1\\2\end{pmatrix}.
\]
Consider \(AB\), \(BA\), \(Au\), \(Bu\), and \(u^{\mathsf T}B\).
\end{enumerate}

\section{Problem Set 2: Determinants}

\begin{enumerate}
\item Compute
\[
\det\begin{pmatrix}4&-3\\2&5\end{pmatrix}.
\]

\item Find the determinant of
\[
A=\begin{pmatrix}
1&2&0\\
-1&3&4\\
2&0&5
\end{pmatrix}
\]
by expanding along the first row.

\item Compute the determinant of
\[
B=\begin{pmatrix}
2&0&0\\
-1&3&0\\
4&5&-2
\end{pmatrix}
\]
without a Laplace expansion. Explain why your method works.

\item Expand
\[
\det\begin{pmatrix}
1&0&2\\
3&-1&4\\
0&5&2
\end{pmatrix}
\]
along the row or column that gives the shortest calculation.

\item A matrix \(B\) is obtained from a \(3\times3\) matrix \(A\) by interchanging two rows. If \(\det A=-7\), find \(\det B\).

\item A matrix \(C\) is obtained from a \(4\times4\) matrix \(A\) by multiplying one row by \(5\). If \(\det A=3\), find \(\det C\).

\item A matrix \(D\) is obtained from \(A\) by adding \(4\) times the first row to the third row. If \(\det A=-2\), determine \(\det D\).

\item Suppose \(\det A=3\) and \(\det B=-4\) for two \(3\times3\) matrices. Find \(\det(AB)\), \(\det(A^{\mathsf T})\), and \(\det(2A)\).

\item Determine whether
\[
A=\begin{pmatrix}
1&2&3\\
2&4&6\\
0&1&1
\end{pmatrix}
\]
is singular without carrying out a full determinant expansion.

\item Find all values of \(k\) for which
\[
A(k)=\begin{pmatrix}
k&2\\
3&k-1
\end{pmatrix}
\]
is singular.

\item Find \(a\) if
\[
\det\begin{pmatrix}
a&1\\
2&3
\end{pmatrix}=10.
\]

\item Compute
\[
\det\begin{pmatrix}
1&2&3\\
0&4&5\\
0&0&6
\end{pmatrix}
\]
and explain the general rule suggested by the result.

\item Let
\[
A=\begin{pmatrix}
1&2&0\\
3&1&4\\
2&-1&5
\end{pmatrix}.
\]
Use row operations to transform \(A\) into an upper triangular matrix, keeping track of every change to the determinant. Then compute \(\det A\).

\item A \(3\times3\) matrix has determinant \(6\). What happens to its determinant if:
\begin{enumerate}
\item the first and second rows are interchanged;
\item the third row is multiplied by \(-2\);
\item twice the first row is added to the second row?
\end{enumerate}

\item The columns of a \(3\times3\) matrix satisfy \(c_3=2c_1-c_2\). What is the determinant? Explain the structural reason.

\item The linear transformation
\[
T(x)=
\begin{pmatrix}
2&1\\
0&3
\end{pmatrix}x
\]
acts on the plane. By what factor does \(T\) scale areas? Does it preserve orientation?

\item Find the area of the parallelogram spanned by
\[
u=(3,1),
\qquad
v=(-2,4)
\]
using a determinant.

\item Compute the Vandermonde determinant
\[
\det\begin{pmatrix}
1&1&1\\
x_1&x_2&x_3\\
x_1^2&x_2^2&x_3^2
\end{pmatrix}
\]
for \(x_1=1\), \(x_2=2\), and \(x_3=4\), first directly and then using the product formula.

\item Consider
\[
A(t)=\begin{pmatrix}
1&t&0\\
0&1&t\\
t&0&1
\end{pmatrix}.
\]
Compute \(\det A(t)\) and determine all real values of \(t\) for which the matrix is singular.

\item For
\[
A=\begin{pmatrix}
1&2&3\\
2&5&7\\
0&1&4
\end{pmatrix},
\]
compute the determinant in two different ways: by a Laplace expansion and by row reduction. Compare the amount of work and explain which method is more efficient.
\end{enumerate}

\section{Problem Set 3: Matrix Inversion}

\begin{enumerate}
\item Decide whether each matrix is invertible:
\[
A=\begin{pmatrix}2&1\\3&4\end{pmatrix},
\qquad
B=\begin{pmatrix}1&2\\2&4\end{pmatrix}.
\]

\item Find the inverse of
\[
A=\begin{pmatrix}3&1\\2&1\end{pmatrix}
\]
using the \(2\times2\) inverse formula.

\item Verify your answer to the previous problem by computing both \(AA^{-1}\) and \(A^{-1}A\).

\item Find the inverse of
\[
D=\operatorname{diag}(2,-3,5).
\]

\item Find all values of \(k\) for which
\[
A(k)=\begin{pmatrix}
k&1\\
2&k
\end{pmatrix}
\]
has an inverse.

\item Use Gauss--Jordan elimination to find the inverse of
\[
A=\begin{pmatrix}
1&2\\
3&5
\end{pmatrix}.
\]

\item Use Gauss--Jordan elimination to find the inverse of
\[
A=\begin{pmatrix}
1&1&0\\
0&1&1\\
1&0&1
\end{pmatrix}.
\]

\item Find the adjugate of
\[
A=\begin{pmatrix}
4&1\\
2&3
\end{pmatrix}
\]
and use it to compute \(A^{-1}\).

\item Let
\[
A=\begin{pmatrix}2&1\\1&1\end{pmatrix},
\qquad
b=\begin{pmatrix}5\\3\end{pmatrix}.
\]
Solve \(Ax=b\) using \(A^{-1}\).

\item Solve
\[
XA=B
\]
for \(X\), where
\[
A=\begin{pmatrix}1&1\\0&2\end{pmatrix},
\qquad
B=\begin{pmatrix}3&1\\4&2\end{pmatrix}.
\]
Pay attention to the side on which \(A^{-1}\) must be multiplied.

\item Solve
\[
2X-A=B
\]
for \(X\), where
\[
A=\begin{pmatrix}1&0\\2&-1\end{pmatrix},
\qquad
B=\begin{pmatrix}3&4\\0&5\end{pmatrix}.
\]

\item If
\[
A^{-1}=\begin{pmatrix}2&-1\\0&3\end{pmatrix}
\quad\text{and}\quad
B^{-1}=\begin{pmatrix}1&0\\4&2\end{pmatrix},
\]
find \((AB)^{-1}\).

\item Prove by direct multiplication that
\[
(A^{\mathsf T})^{-1}=(A^{-1})^{\mathsf T}
\]
for
\[
A=\begin{pmatrix}1&2\\3&5\end{pmatrix}.
\]

\item Find the inverse of the elementary matrix
\[
E=\begin{pmatrix}
1&0&0\\
0&1&0\\
-4&0&1
\end{pmatrix}
\]
and explain which row operation \(E\) represents.

\item A proposed inverse of
\[
A=\begin{pmatrix}2&-1\\1&0\end{pmatrix}
\]
is
\[
C=\begin{pmatrix}0&1\\-1&2\end{pmatrix}.
\]
Check whether the proposal is correct.

\item Explain why a matrix with two identical rows cannot have an inverse. Use both a determinant argument and an information-loss interpretation.

\item The transformation
\[
y=
\begin{pmatrix}2&1\\1&1\end{pmatrix}x
\]
maps an unknown vector \(x\) to
\[
y=\begin{pmatrix}8\\5\end{pmatrix}.
\]
Recover \(x\).

\item Choose the most efficient method for finding the inverse of each matrix and justify the choice:
\[
A=\begin{pmatrix}5&0\\0&-2\end{pmatrix},
\quad
B=\begin{pmatrix}1&2\\3&7\end{pmatrix},
\quad
C=\begin{pmatrix}1&1&0\\0&1&1\\0&0&2\end{pmatrix}.
\]
Then compute the inverses.

\item Let
\[
A=\begin{pmatrix}
1&a\\
2&2a
\end{pmatrix}.
\]
Explain, without attempting Gauss--Jordan elimination, why \(A^{-1}\) does not exist for any real \(a\).

\item For
\[
A=\begin{pmatrix}
1&2\\
3&4
\end{pmatrix},
\]
find \(A^{-1}\), solve \(Ax=b\) for
\[
b=\begin{pmatrix}5\\11\end{pmatrix},
\]
and verify the solution in the original equation. Finally, explain what the inverse accomplishes conceptually.
\end{enumerate}

\section{Problem Set 4: Systems of Linear Equations}

\begin{enumerate}
\item Solve by substitution:
\[
\begin{cases}
x+y=7,\\
2x-y=2.
\end{cases}
\]

\item Write the augmented matrix for
\[
\begin{cases}
2x-y+3z=5,\\
x+4y-z=2,\\
3x+2z=7.
\end{cases}
\]

\item Solve by Gaussian elimination:
\[
\begin{cases}
x+2y=5,\\
3x-y=4.
\end{cases}
\]

\item Solve by Gaussian elimination:
\[
\begin{cases}
x+y+z=6,\\
2x-y+z=3,\\
x+3y-2z=1.
\end{cases}
\]

\item Reduce the augmented matrix
\[
\left(
\begin{array}{cc|c}
1&2&3\\
2&4&8
\end{array}
\right)
\]
to echelon form and classify the system.

\item Reduce
\[
\left(
\begin{array}{ccc|c}
1&2&-1&3\\
2&4&-2&6\\
0&1&1&2
\end{array}
\right)
\]
and express all solutions using a parameter.

\item Determine whether
\[
\begin{cases}
x-2y=1,\\
2x-4y=5
\end{cases}
\]
has no solution, one solution, or infinitely many solutions. Interpret the result geometrically.

\item Solve
\[
\begin{cases}
2x+y=4,\\
x-3y=-5
\end{cases}
\]
using Cramer's rule.

\item Solve the system in the previous problem using the inverse matrix method. Compare the two methods.

\item For which values of \(k\) does
\[
\begin{cases}
kx+y=2,\\
2x+ky=3
\end{cases}
\]
have a unique solution?

\item Classify the system according to the value of \(a\):
\[
\begin{cases}
x+y=2,\\
2x+2y=a.
\end{cases}
\]

\item Find all values of \(k\) for which
\[
\begin{cases}
x+y+z=1,\\
2x+2y+2z=2,\\
x+y+kz=3
\end{cases}
\]
is consistent. Describe the solution set in each consistent case.

\item The row-reduced augmented matrix of a system is
\[
\left(
\begin{array}{cccc|c}
1&0&2&0&3\\
0&1&-1&0&4\\
0&0&0&1&2
\end{array}
\right).
\]
Identify the pivot and free variables and write the complete solution parametrically.

\item Determine whether the point \((2,-1)\) is the solution of
\[
\begin{cases}
3x+y=5,\\
x-2y=4.
\end{cases}
\]
If it is not, find the correct solution.

\item Find the intersection point of the lines
\[
2x-y-1=0,
\qquad
x+3y-9=0.
\]

\item Find the intersection of the three planes
\[
\begin{cases}
x+y+z=3,\\
x-y+z=1,\\
2x+z=2.
\end{cases}
\]

\item A cinema sells adult tickets for \(12\) units and student tickets for \(8\) units. In one evening it sells \(180\) tickets for a total of \(1840\) units. Form and solve a linear system to find the number of each type of ticket.

\item A system reduces to a row
\[
\begin{pmatrix}0&0&0\mid 6\end{pmatrix}.
\]
What does this row mean, and what can be concluded about the original system?

\item Construct a system of two linear equations in two variables that has:
\begin{enumerate}
\item exactly one solution;
\item no solution;
\item infinitely many solutions.
\end{enumerate}
Explain how the coefficients produce each case.

\item Consider
\[
\begin{cases}
x+ay=1,\\
ax+y=1.
\end{cases}
\]
Classify the system for every real value of \(a\). For each case, determine whether the solution is unique, nonexistent, or non-unique, and find the solution set where possible.
\end{enumerate}
